% ==============================================================
% Das Fragedifferential (DFD) -- Version 2.2
% Ontologisches Individuationsprinzip des Ontophänomenalismus
% Kompiliert mit: pdfLaTeX
% ==============================================================

\documentclass[11pt,a4paper]{article}

% ---------- Grundpakete ----------
\usepackage[utf8]{inputenc}
\usepackage[T1]{fontenc}
\usepackage[german]{babel}
\usepackage{amsmath,amssymb,amsthm}
\usepackage{geometry}
\usepackage{parskip}
\usepackage{enumitem}
\usepackage{booktabs}
\usepackage{xcolor}
\usepackage{hyperref}

% ---------- Layout ----------
\geometry{margin=2.5cm}
\setlist{itemsep=0.3em, topsep=0.3em}
\hypersetup{
    colorlinks=true,
    linkcolor=blue,
    urlcolor=blue,
    pdftitle={Das Fragedifferential (DFD) -- Version 2.2},
    pdfauthor={Forschungsprogramm Ontoph\"{a}nomenalismus}
}

% ---------- Theorem-Umgebungen ----------
\theoremstyle{plain}
\newtheorem{theorem}{Theorem}[section]
\newtheorem{proposition}{Proposition}[section]
\newtheorem{corollary}{Korollar}[section]

\theoremstyle{definition}
\newtheorem{definition}{Definition}[section]
\newtheorem{remark}{Bemerkung}[section]
\newtheorem{methodennote}{Methodennote}[section]

% ---------- Titel ----------
\title{
  \textbf{Das Fragedifferential (DFD)}\\[0.6em]
  \large Das ontologische Individuationsprinzip des Ontoph\"{a}nomenalismus\\[0.4em]
  {\normalsize Version~2.2}
}
\author{
  \textbf{Forschungsprogramm Ontoph\"{a}nomenalismus}\\
  \small Siehe TRANSPARENCY f\"{u}r methodologische Urheberschaft und KI-Einsatz.
}
\date{Juni 2026}

\begin{document}

\maketitle

\begin{abstract}
Dieses Dokument etabliert das Fragedifferential (DFD) als eigenst\"{a}ndiges Modul
innerhalb der Architektur des Ontoph\"{a}nomenalismus (OP). Es l\"{o}st das DFD aus
dem begrifflichen Schatten des Explanatory Gap (EG-1) heraus und begr\"{u}ndet es als
das prim\"{a}re informationelle Individuationsprinzip des Systems.

Das DFD beantwortet die inverse metaphysische Herausforderung des OP: Wie differenziert
sich innerhalb eines raumzeitlosen, kontinuierlichen und unpers\"{o}nlichen Gesamtfeldes
\"{u}berhaupt eine isolierte, begrenzte und individuelle Perspektive heraus? Individuation
wird dabei nicht als Hinzuf\"{u}gung einer substanziellen Eigenschaft verstanden, sondern
als konstitutive, strukturierte Differenz zum Absoluten.

Es gelten die methodologischen Konventionen gem\"{a}\ss{} MKO --- insbesondere die
Richtlinie zur Vermeidung vorzeitiger Kontamination: Das Fragedifferential beschreibt
die informationelle Differenz zwischen der lokalen Zustandsverteilung $Q_E$ einer
projizierten Entit\"{a}t (entstanden durch DPR) und dem Gesamtfeld $P$ --- es operiert
damit auf der Ebene nach der Projektion, nicht im ontisch metrikfreien OPP selbst.
\end{abstract}

\tableofcontents
\newpage

%==============================================================
\section{Motivation: Das fundamentale Individuationsproblem}
%==============================================================

Ein radikal monistischer Ansatz wie der OP, der von einem ungeteilten primordialen
Denkfeld als absolutem Urgrund ausgeht, steht vor einer inversen metaphysischen
Herausforderung gegen\"{u}ber dem klassischen Materialismus. W\"{a}hrend der
Physikalismus erkl\"{a}ren muss, wie aus toter Materie subjektives Erleben entsteht,
muss der OP das exakte Gegenteil herleiten: Wie bildet sich innerhalb eines
raumzeitlosen, kontinuierlichen und unpers\"{o}nlichen Gesamtfeldes \"{u}berhaupt eine
isolierte, begrenzte und individuelle Perspektive heraus?

Das Fragedifferential ($\mathcal{D}(E)$) liefert den konzeptuellen Rahmen f\"{u}r diese
Frage. Es k\"{o}nnte begr\"{u}nden, warum die Differenzierung in scheinbar getrennte
Perspektiven strukturell notwendig ist, sobald das Feld sich selbst aus einer
lokalisierten, projizierten Geometrie heraus reflektiert. Ohne ein dediziertes
Individuationsprinzip bliebe jeder Monismus eine konturlose Identit\"{a}tsphilosophie,
die unf\"{a}hig w\"{a}re, die ph\"{a}nomenale Vielheit der erlebten Welt zu erkl\"{a}ren.

%==============================================================
\section{Historische und systematische Einordnung}
%==============================================================

Um die Eigenst\"{a}ndigkeit des DFD zu verdeutlichen, wird es in den historischen Kontext
der Philosophie und der modernen Kognitionswissenschaft gestellt:

\begin{itemize}
  \item \textbf{Gottfried Wilhelm Leibniz:} Seine \emph{Monadenlehre} l\"{o}st die
        Individuation \"{u}ber die Identit\"{a}t des Ununterscheidbaren. Die Monade
        spiegelt das Universum aus einer spezifischen Perspektive. Das DFD formalisiert
        und dynamisiert diesen Gedanken informationstheoretisch: Die Perspektive
        \emph{k\"{o}nnte} die quantifizierbare informationelle Abweichung vom Ganzen
        sein.

  \item \textbf{Edmund Husserl:} Die ph\"{a}nomenologische Intentionalit\"{a}t setzt
        immer ein Ich-Zentrum voraus. Das DFD k\"{o}nnte die ontologische Basis f\"{u}r
        diese Intentionalit\"{a}t liefern, indem es zeigt, dass das
        \glqq{}Gerichtet-Sein\grqq{} auf ein Objekt erst durch ein strukturelles
        Informationsdefizit gegen\"{u}ber dem unreduzierten Phasenraum entsteht.

  \item \textbf{Alfred North Whitehead:} In seiner Prozessphilosophie konstituieren
        sich \glqq{}wirkliche Einzelwesen\grqq{} durch Selektion und Abgrenzung
        (Pr\"{a}hension) von Daten. Das DFD \"{u}bersetzt diesen Gedanken in die Sprache
        einer informationstheoretischen Geometrie --- als Analogie (MKO, Ebene~4),
        nicht als Identifikation.

  \item \textbf{Integrierte Informationstheorie (IIT, Tononi):} W\"{a}hrend die IIT
        Bewusstsein \"{u}ber die integrierte Information ($\Phi$) innerhalb eines
        Systems misst, geht das DFD den umgekehrten Weg: Es misst nicht, wie viel
        Information im System integriert ist, sondern wie viel Information dem System
        im Vergleich zum unreduzierten Urgrund \emph{fehlt}. Das DFD k\"{o}nnte in
        diesem Sinne als \glqq{}Exformationstheorie\grqq{} der Individuation verstanden
        werden --- eine Analogie (MKO, Ebene~4), die der formalen Ausarbeitung durch
        die PST bedarf.
\end{itemize}

%==============================================================
\section{Ontologische Definition des DFD}
%==============================================================

Das Fragedifferential ist im OP \textbf{keine} rein epistemische Kategorie. Es beschreibt
nicht das blo\ss{}e psychologische Unwissen eines bereits existierenden Subjekts. Vielmehr
ist das DFD ein \emph{ontologisches Konstitutionsprinzip}: Die informationelle L\"{u}cke
ist nicht ein Defizit eines vorhandenen Wesens, sondern die strukturelle Bedingung, unter
der Individualit\"{a}t \"{u}berhaupt erst entsteht.

Eine individuelle Entit\"{a}t existiert nicht unabh\"{a}ngig und \emph{hat} dann offene
Fragen; die spezifische, strukturierte Anordnung ihrer informationellen Differenz zum
Gesamtfeld \textbf{ist} ihre ontologische Grenze. Individualit\"{a}t im OP bedeutet:
ein lokalisierter Fokus im projizierten Feld zu sein, der durch eine pr\"{a}zise
definierte informationelle Asymmetrie gegen\"{u}ber dem Ganzen konstituiert wird.

%==============================================================
\section{Mathematische Heuristik}
%==============================================================

\begin{methodennote}[Methodischer Vorbehalt --- gem\"{a}\ss{} MKO, Ebene~3]
Das Fragedifferential ist prim\"{a}r ein ontologischer Begriff des OP. Die nachfolgende
mathematische Formulierung mittels der Kullback-Leibler-Divergenz stellt keine Behauptung
dar, dass beide Begriffe ontologisch identisch sind. Die KL-Divergenz dient als erste
informationstheoretische Heuristik --- als formales Modell im Sinne von MKO. Ob die
PST dieselbe Gr\"{o}\ss{}e, eine Verallgemeinerung oder eine grundlegend andere
mathematische Struktur verwenden wird, bleibt eine offene Forschungsfrage.
\end{methodennote}

\begin{definition}[Das Fragedifferential]
Sei $\mathcal{M}_{\mathrm{OPP}}$ der ontoph\"{a}nomenale Phasenraum und $P$ die
fundamentale Zustandsverteilung des ungeteilten primordialen Denkfeldes. Sei ferner
$Q_E$ die durch die Projektive Reduktion (DPR) entstehende lokale Zustandsverteilung
einer individuierten Entit\"{a}t $E$.

Als erste mathematische Heuristik definieren wir das Fragedifferential $\mathcal{D}(E)$
durch die Kullback-Leibler-Divergenz:
\[
  \mathcal{D}(E) \;:=\; D_{\mathrm{KL}}(Q_E \,\|\, P)
  \;=\; \sum_{x \in \mathcal{X}} Q_E(x) \log \left( \frac{Q_E(x)}{P(x)} \right)
\]
\end{definition}

\begin{remark}[Informationstheoretische Interpretation]
Die Kullback-Leibler-Divergenz ist nicht symmetrisch
($D_{\mathrm{KL}}(Q_E \| P) \neq D_{\mathrm{KL}}(P \| Q_E)$). Diese Richtung ist
bewusst gew\"{a}hlt: Sie beschreibt den informationellen Abstand der lokalen Perspektive
$Q_E$ gegen\"{u}ber dem vollst\"{a}ndigen Denkfeld $P$. Das DFD k\"{o}nnte daher als
quantitatives Ma\ss{} der \emph{konstitutiven Offenheit} einer Perspektive interpretiert
werden. Die bildhafte Formulierung der \glqq{}offenen Fragen\grqq{} ist keine
mathematische Definition, sondern eine philosophische Interpretation dieser strukturellen
Differenz (MKO, Ebene~2).
\end{remark}

%==============================================================
\section{Eigenschaften des Fragedifferentials}
%==============================================================

\subsection{Strikt positive Asymmetrie}

Die unumkehrbare Gerichtetheit der Perspektive spiegelt sich in der mathematischen
Asymmetrie wider: Eine begrenzte Entit\"{a}t kann das Ganze nur unter permanentem
Informationsverlust approximieren, w\"{a}hrend das Ganze die Entit\"{a}t stets
vollst\"{a}ndig umfasst.

\subsection{Die zwei fundamentalen Grenzf\"{a}lle}

Aus der mathematischen Formulierung ergeben sich zwei extreme konzeptuelle Zust\"{a}nde:

\begin{enumerate}
  \item \textbf{Vollst\"{a}ndige Identit\"{a}t ($\mathcal{D}(E) = 0$):} Wenn die
        informationelle Differenz vollst\"{a}ndig verschwindet, kollabiert $Q_E$ in $P$.
        Das bedeutet das Ende der Individuation. Dieser Grenzfall ist strukturell
        ausgeschlossen f\"{u}r jede endliche Perspektive (vgl.\ EG-1, MKO).

  \item \textbf{Maximale Isolation ($\mathcal{D}(E) \to \infty$):} Je gr\"{o}\ss{}er
        die Divergenz, desto restriktiver der projektive Filter. Die Entit\"{a}t
        k\"{o}nnte sich als vollkommen isoliertes System in einer ihr fremden Umwelt
        erfahren. Dieser Grenzfall ist konzeptuell, nicht empirisch behauptet
        (MKO, Ebene~2).
\end{enumerate}

%==============================================================
\section{Das Prinzip der Individuation}
%==============================================================

\begin{proposition}[Kriterium der Individuation]
Eine Entit\"{a}t $E$ existiert genau dann als individuiertes System innerhalb des
Urgrundes, wenn gilt:
\[
  \mathcal{D}(E) \;>\; 0
\]
\end{proposition}

\begin{remark}[Konsequenz f\"{u}r den Monismus]
Hieraus folgt: Begrenzung und strukturelle Unvollst\"{a}ndigkeit sind keine M\"{a}ngel
des Seins, sondern die notwendigen Bedingungen f\"{u}r die Existenz von Individualit\"{a}t
und Perspektive. Der OP erkl\"{a}rt damit die ph\"{a}nomenale Vielheit nicht trotz,
sondern durch seinen Monismus.
\end{remark}

%==============================================================
\section{Erkenntnistheoretische Konsequenzen: Die Begriffstriade}
%==============================================================

Das DFD sch\"{u}tzt den OP vor dem Abgleiten in einen naiven Relativismus, ohne in
einen dogmatischen Absolutismus zu verfallen (vgl.\ EG-1, Sektion~7):

\begin{table}[h!]
\centering
\small
\begin{tabular}{p{2.8cm}p{4.5cm}p{5.5cm}}
  \toprule
  \textbf{Kategorie} & \textbf{Informationeller Zustand} & \textbf{Ontologische
  Entsprechung} \\
  \midrule
  \textbf{Wahrheit} & Unreduziertes Gesamtfeld ($P$) &
    Das zeitlose, kontinuierliche OPP. Invarianter Bezugspunkt aller Perspektiven. \\[4pt]
  \textbf{Objektivit\"{a}t} & Grenzfall $\mathcal{D}(E) = 0$ &
    Asymptotisches Ideal; w\"{u}rde Kollaps der Perspektive bedeuten. \\[4pt]
  \textbf{Perspektivit\"{a}t} & $\mathcal{D}(E) > 0$ &
    Notwendige Bedingung f\"{u}r lokalisierte, individuelle Erkenntnis. \\
  \bottomrule
\end{tabular}
\caption{Die ontoph\"{a}nomenalistische Begriffstriade.}
\end{table}

%==============================================================
\section{Integration in die Modularchitektur}
%==============================================================

Das DFD ist ein Interface-Objekt, das zwischen mehreren Modulen vermittelt:

\begin{itemize}
  \item \textbf{OAD (Ontologische Attraktoren-Dynamik):} Die OAD liefert die stabilen
        Konfigurationen hoher struktureller Koh\"{a}renz im OPP, deren projektive
        Manifestationen das DFD charakterisiert.
  \item \textbf{DPR (Die Projektive Reduktion):} Die lokale Zustandsverteilung $Q_E$
        entsteht durch den Projektionsoperator der DPR. DFD und DPR sind damit eng
        verkn\"{u}pft: Die Projektion erzeugt die Differenz, die das DFD misst.
  \item \textbf{DPD (Die Ph\"{a}nomenale Differenz):} Das DFD liefert das
        ontologische Fundament f\"{u}r die Ph\"{a}nomenale Differenz: Das Gef\"{a}lle
        zwischen zwei Attraktoren ist eine spezifische Form der informationellen
        Asymmetrie, die das DFD allgemein beschreibt.
  \item \textbf{EG-1 (Der Explanatory Gap):} Das DFD liefert dort das
        Individuationsprinzip und begr\"{u}ndet, warum der Explanatory Gap keine
        schlie\ss{}bare L\"{u}cke ist, sondern die konstitutive Grenze jeder endlichen
        Perspektive.
  \item \textbf{PST (Projektive Strukturtheorie):} Die PST \"{u}bernimmt die Aufgabe,
        $\mathcal{D}(E)$ mathematisch pr\"{a}zise zu definieren und aus dem OPP
        herzuleiten.
\end{itemize}

%==============================================================
\section{Physikalische Analogien und Hypothesen}
%==============================================================

\begin{methodennote}[Status dieser Sektion --- gem\"{a}\ss{} MKO, Ebene~4 und~5]
Die folgenden \"{U}berlegungen sind Analogien und Hypothesen im Sinne von MKO. Sie
behaupten keine physikalischen Tatsachen, sondern skizzieren m\"{o}gliche
Anschlusspunkte f\"{u}r die PST.
\end{methodennote}

\begin{itemize}
  \item \textbf{Analogie: Subjektiver Zeitpfeil (MKO, Ebene~4).} Warum vergeht Zeit
        f\"{u}r ein Individuum linear? Eine m\"{o}gliche Interpretation: Weil das System
        durch seine Begrenzung gezwungen sein k\"{o}nnte, die raumzeitlose
        Zustandsverteilung des OPP sequenziell zu approximieren --- \emph{analog zu}
        einer asymptotischen Ann\"{a}herung an $P$.

  \item \textbf{Hypothese: Lernen und lokale Informationsgewinnung (MKO, Ebene~5).}
        Falls der OP zutrifft, w\"{a}re zu erwarten, dass jeder Akt des Lernens das
        Fragedifferential lokal verringert ($\Delta \mathcal{D}(E) < 0$). Ob und wie
        das empirisch pr\"{u}fbar ist, bleibt eine offene Forschungsfrage der PST.
\end{itemize}

%==============================================================
\section{Offene Forschungsfragen f\"{u}r die PST}
%==============================================================

\begin{enumerate}
  \item Welche verallgemeinerte informationstheoretische Metrik --- etwa
        Von-Neumann-Entropie auf Operatoralgebren oder eine andere raumzeitlose Struktur
        --- beschreibt das Feld pr\"{a}zise?
  \item Wie koppelt die Dynamik des DFD an die Projektionsoperatoren der DPR?
  \item L\"{a}sst sich ein Informationsgewinn einer Perspektive als messbare
        Zustands\"{a}nderung in einem raumzeitlichen Quantensystem operationalisieren?
  \item Wie verh\"{a}lt sich das DFD zu $\Phi$ der IIT --- Komplement, Verallgemeinerung
        oder grundlegend andere Kategorie?
\end{enumerate}

\end{document}
